\begin{definition}[Sub-protocol specifications]
  \label{def:aba-subspecs}
  \lean{PLTS.ABA.WCC.Step}\lean{PLTS.ABA.WCC.specFamily}\lean{PLTS.ABA.GBCA.Step}\lean{PLTS.ABA.GBCA.specFamily}\leanok
  \uses{def:aba-labels, def:idle-family}
  Transition Systems 2 and 3, as round-indexed instances and the $\mathbb{N}$-indexed families of Definition~\ref{def:idle-family} over them: GBCA with internal binding and the graded returns $A\,b$, $B\,b$ and $C$, WCC with the $\varepsilon$-probabilistic coin resolution.
  The GBCA side is an LTS.\lean{PLTS.ABA.GBCA.specFamily_isLTS}\leanok{}
  A GBCA instance binds by exclusion (D19): its state carries the set $dead \subseteq \{0,1\}$ of bits the instance can no longer hand out, in place of the source's bound value, written only by an internal kill enabled from $dead = \emptyset$ alone, so the kill fires at most once per instance; beside it the state carries the pending call $call$ and the return flag $ret$ of each process, the grade lock $grade$, latched by an $A$-return and by a $C$-return exclusively, and the corrupted set.
  That kill, the dissenting bit at a $B$-return and each of the two bits at a $C$-return carry an $f{+}1$ $F$-blind support guard, $\mathrm{SuppOK}$: $f{+}1$ distinct callers-or-corrupted of the bit in question (D14, D15).
  Full rule tables: \leandecl{PLTS.ABA.GBCA.specFamily} and \leandecl{PLTS.ABA.WCC.specFamily}.
\end{definition}
